\documentclass[12pt,a4paper]{article}

% --- Paquetes ---
\usepackage[utf8]{inputenc}
\usepackage[T1]{fontenc}
\usepackage[spanish,provide=*]{babel}
\usepackage{amsmath, amssymb, amsthm}
\usepackage{geometry}
\usepackage{hyperref}
\usepackage{enumitem}
\usepackage{setspace}
\usepackage{fancyhdr}
\usepackage{booktabs}
\usepackage{xcolor}
\usepackage{tcolorbox}
\usepackage{tikz}
\usetikzlibrary{arrows.meta, positioning, shapes}

\geometry{margin=2.5cm}
\onehalfspacing

% --- Teoremas y entornos ---
\newtheorem{theorem}{Teorema}
\newtheorem{lemma}[theorem]{Lema}
\newtheorem{corollary}[theorem]{Corolario}
\newtheorem{proposition}[theorem]{Proposición}
\theoremstyle{definition}
\newtheorem{definition}{Definición}
\newtheorem{example}{Ejemplo}
\theoremstyle{remark}
\newtheorem{remark}{Observación}

% --- Cabecera ---
\pagestyle{fancy}
\fancyhf{}
\rhead{NP-Completitud}
\lhead{Problema de la 3-Partición}
\rfoot{\thepage}

% --- Comandos útiles ---
\newcommand{\TPART}{\textsc{3-Partition}}
\newcommand{\FPART}{\textsc{4-Partition}}
\newcommand{\DM}{\textsc{3DM}}
\newcommand{\PART}{\textsc{Partition}}
\newcommand{\SUMA}{\textsc{Subset-Sum}}

% Caja para "intuición" (explicación para el lector)
\newtcolorbox{intuicion}[1][]{
  colback=yellow!8, colframe=yellow!55!black, coltitle=black,
  fonttitle=\bfseries, title={Intuición}, #1}

% ============================================================
\begin{document}
% ============================================================

\begin{titlepage}
  \centering
  \vspace*{2cm}
  {\LARGE\bfseries Demostración de NP-Completitud del\\[0.5em]
  Problema de la 3-Partición\par}
  \vspace{1.5cm}
  {\large Trabajo de Complejidad Computacional\par}
  \vspace{0.5cm}
  {\large Ismael Sallami Moreno\par}
  \vspace{0.5cm}
  {\large Universidad de Granada\par}
  \vspace{2cm}
  {\normalsize 2026\par}
  \vfill
  \begin{tcolorbox}[colback=gray!10, colframe=gray!50, width=0.85\textwidth]
  \small\centering
  \textbf{Referencia principal:} M.\,R. Garey y D.\,S. Johnson,
  \textit{Computers and Intractability: A Guide to the Theory of NP-Completeness},
  W.H. Freeman, 1979 (Teoremas 4.3 y 4.4, pp.~96--101; problema [SP15]).
  \end{tcolorbox}
\end{titlepage}

\tableofcontents
\newpage

% ============================================================
\section{Introducción y Motivación}
% ============================================================

El \textbf{Problema de la 3-Partición} (\TPART) es uno de los problemas
NP-completos más útiles de la teoría de la complejidad computacional.
A diferencia de otros problemas de partición de enteros, como el Problema de la
Mochila o el Problema de la Partición clásica (en dos partes), la 3-Partición
posee una propiedad más fuerte: es \textbf{fuertemente NP-completo}
(\emph{strongly NP-complete}).

\begin{intuicion}
Casi todos los problemas NP-completos sobre números (Mochila, Partición en dos)
se vuelven fáciles si los números de la entrada son pequeños: existe un algoritmo
de \emph{programación dinámica} cuyo coste depende del valor de los números, no
solo de cuántos hay. Cuando los números son pequeños, ese algoritmo es rápido.
La 3-Partición es especial precisamente porque \textbf{sigue siendo difícil
aunque los números sean pequeños}. Esto es lo que significa ``fuertemente''
NP-completo, y es lo que la convierte en una herramienta tan buena para demostrar
que \emph{otros} problemas son difíciles.
\end{intuicion}

Esta distinción tiene importancia teórica y práctica:
\begin{itemize}
  \item Los problemas NP-completos ``débiles'' (como Mochila o
        \PART\ clásica) admiten algoritmos de programación dinámica
        en tiempo \emph{pseudopolinómico}: son tratables cuando
        los números son pequeños (acotados polinómicamente en el tamaño de la
        entrada).
  \item \TPART\ es intratable \emph{incluso cuando los números son pequeños}.
        La programación dinámica no ofrece ninguna ventaja, salvo que
        $\mathbf{P}=\mathbf{NP}$.
\end{itemize}

El problema fue catalogado y demostrado fuertemente NP-completo por
Garey y Johnson, quienes lo usaron como pieza intermedia para analizar
problemas de planificación multiprocesador. En su libro de 1979 aparece como
problema [SP15], y la demostración ocupa los Teoremas 4.3 y 4.4
(pp.~96--101), que es la fuente que sigue este documento.

En este trabajo se presenta la demostración formal siguiendo \emph{exactamente}
la cadena de reducciones de Garey y Johnson:
\[
  \DM \;\leq_p\; \FPART \;\leq_p\; \TPART,
\]
donde \DM\ es el Emparejamiento Tridimensional (3-Dimensional Matching).

\begin{remark}
La cadena pasa por un problema intermedio, \FPART\ (la 4-Partición). Una pregunta
natural es \emph{por qué hacen falta dos reducciones y no una sola}. La respuesta
se discute en detalle en la Sección~\ref{sec:dos-reducciones}, pero la idea breve
es: en principio bastaría una sola reducción desde cualquier problema NP-completo;
sin embargo, construir directamente $\DM \leq_p \TPART$ es técnicamente muy
complicado, y el escalón intermedio \FPART\ hace cada paso manejable. Por la
transitividad de $\leq_p$, dos reducciones encadenadas equivalen a una.
\end{remark}

% ============================================================
\section{Preliminares}
% ============================================================

\subsection{Definiciones de NP-completitud}

\begin{definition}[Clase NP]
Un problema de decisión $\Pi$ pertenece a la clase \textbf{NP} si existe
una máquina de Turing no determinista que lo resuelve en tiempo polinómico.
De forma equivalente, dado un \emph{certificado} (solución candidata), existe un
algoritmo determinista que lo \emph{verifica} en tiempo polinómico en el tamaño
de la entrada.
\end{definition}

\begin{definition}[Reducción polinómica]
Un problema $\Pi_1$ se \emph{reduce polinómicamente} a $\Pi_2$
($\Pi_1 \leq_p \Pi_2$) si existe una función $f$ computable en tiempo
polinómico tal que:
\[
  x \in \Pi_1 \iff f(x) \in \Pi_2.
\]
\end{definition}

\begin{intuicion}
Una reducción $\Pi_1 \leq_p \Pi_2$ es una ``traducción'' eficiente: convierte
cada pregunta de $\Pi_1$ en una pregunta equivalente de $\Pi_2$, de modo que la
respuesta (sí/no) se conserva. Si supiéramos resolver $\Pi_2$ rápido, resolveríamos
$\Pi_1$ rápido. Por eso, si $\Pi_1$ ya es ``difícil'', $\Pi_2$ no puede ser más
fácil. La dirección importa: para probar que $\Pi_2$ es difícil, traducimos un
problema difícil conocido \emph{hacia} $\Pi_2$, nunca al revés.
\end{intuicion}

\begin{definition}[NP-completitud]
$\Pi$ es \textbf{NP-completo} si:
\begin{enumerate}
  \item $\Pi \in \mathbf{NP}$, y
  \item todo problema en NP se reduce polinómicamente a $\Pi$
        (equivalentemente, algún problema NP-completo conocido se reduce a $\Pi$).
\end{enumerate}
\end{definition}

\subsection{NP-completitud fuerte}

\begin{definition}[Número máximo y tamaño de instancia]
Para un problema con números en la entrada, sea $\mathrm{Max}[I]$ el mayor valor
numérico que aparece en la instancia $I$, y $\mathrm{Length}[I]$ la longitud de
la codificación (razonable, en binario) de $I$.
\end{definition}

\begin{definition}[Fuertemente NP-completo]
Un problema $\Pi$ es \textbf{fuertemente NP-completo} si sigue siendo NP-completo
cuando lo restringimos a instancias en las que $\mathrm{Max}[I]$ está acotado por
un polinomio en $\mathrm{Length}[I]$. Equivalentemente: $\Pi$ permanece NP-completo
aunque los números se escriban en \emph{unario}.
\end{definition}

\begin{remark}
La NP-completitud fuerte implica que, salvo que $\mathbf{P}=\mathbf{NP}$,
el problema \textbf{no} admite algoritmo pseudopolinómico (uno cuyo tiempo sea
polinómico en $\mathrm{Length}[I]$ y en $\mathrm{Max}[I]$). Esto lo distingue
radicalmente de \PART\ clásica o Mochila, que sí lo admiten.
\end{remark}

\begin{definition}[Transformación pseudopolinómica]
\label{def:pseudopoli}
Una \emph{transformación pseudopolinómica} de $\Pi$ a $\Pi'$ es una función
$f$ que cumple, además de preservar las respuestas
($I \in Y_\Pi \iff f(I) \in Y_{\Pi'}$):
\begin{enumerate}[label=(\alph*)]
  \item $f$ es computable en tiempo polinómico en las dos variables
        $\mathrm{Max}[I]$ y $\mathrm{Length}[I]$;
  \item existe un polinomio $q_1$ con $q_1(\mathrm{Length}'[f(I)]) \geq \mathrm{Length}[I]$;
  \item existe un polinomio $q_2$ con
        $\mathrm{Max}'[f(I)] \leq q_2(\mathrm{Max}[I], \mathrm{Length}[I])$.
\end{enumerate}
\end{definition}

\begin{lemma}[Garey--Johnson, Lema 4.1]
\label{lem:pseudo}
Si $\Pi$ es fuertemente NP-completo, $\Pi' \in \mathbf{NP}$, y existe una
transformación pseudopolinómica de $\Pi$ a $\Pi'$, entonces $\Pi'$ es
fuertemente NP-completo.
\end{lemma}

\begin{intuicion}
El Lema~\ref{lem:pseudo} es la herramienta que nos permite ``propagar'' la
dureza fuerte a lo largo de la cadena. Una reducción polinómica normal podría
crear números gigantescos (exponenciales), y eso arruinaría la dureza
\emph{fuerte}. La condición (c) de la Definición~\ref{def:pseudopoli} es la clave:
exige que los números que genera la reducción no exploten, sino que queden
acotados polinómicamente. Si cada paso de la cadena respeta esto, la dureza
fuerte sobrevive de principio a fin.
\end{intuicion}

% ============================================================
\section{Definición Formal del Problema}
% ============================================================

\begin{tcolorbox}[colback=blue!5, colframe=blue!40, title={\TPART\ (Problema de la 3-Partición) [SP15]}]
\textbf{Entrada:} Un conjunto $A$ de $3m$ elementos, una cota $B \in \mathbb{Z}^+$,
y un ``tamaño'' $s(a) \in \mathbb{Z}^+$ para cada $a \in A$, tales que:
\begin{enumerate}
  \item $B/4 < s(a) < B/2$ para todo $a \in A$ \quad (condición de acotación), y
  \item $\displaystyle\sum_{a \in A} s(a) = m \cdot B$ \quad (condición de suma).
\end{enumerate}
\textbf{Pregunta:} ¿Puede particionarse $A$ en $m$ conjuntos disjuntos
$S_1, S_2, \ldots, S_m$ tales que, para cada $i$,
$\displaystyle\sum_{a \in S_i} s(a) = B$?
\end{tcolorbox}

\begin{remark}
La condición de acotación $B/4 < s(a) < B/2$ no es decorativa: \textbf{fuerza}
que cada $S_i$ tenga exactamente 3 elementos. En efecto, si todos los tamaños
están estrictamente entre $B/4$ y $B/2$, entonces dos elementos suman menos de
$B$ (a lo sumo casi $B$) y cuatro elementos suman más de $B$; la única forma de
alcanzar \emph{exactamente} $B$ es con tres. Por eso ni siquiera hace falta pedir
$|S_i| = 3$ en el enunciado: se deduce de las cotas.
\end{remark}

\begin{example}[Instancia pequeña correcta]
Sea $m = 2$, $B = 24$ y
$A = \{7, 8, 9, 10, 7, 7\}$, con todos los tamaños estrictamente en
$(B/4, B/2) = (6, 12)$.
Una solución es:
\[
  S_1 = \{7, 8, 9\}, \qquad S_2 = \{10, 7, 7\}, \qquad
  \textstyle\sum S_1 = \sum S_2 = 24 = B. \;\checkmark
\]
Obsérvese que $\sum_{a\in A} s(a) = 48 = 2 \cdot 24 = m B$, como exige la
condición de suma.
\end{example}

% ============================================================
\section{Problemas Auxiliares}
% ============================================================

\subsection{Emparejamiento Tridimensional (\DM)}

\begin{tcolorbox}[colback=green!5, colframe=green!40, title={\DM\ (3-Dimensional Matching) [SP1]}]
\textbf{Entrada:} Tres conjuntos disjuntos $W$, $X$, $Y$ con $|W|=|X|=|Y|=q$,
y un conjunto de tripletas $M \subseteq W \times X \times Y$.\\
\textbf{Pregunta:} ¿Existe un subconjunto $M' \subseteq M$ con $|M'| = q$
tal que cada elemento de $W \cup X \cup Y$ aparece en exactamente una
tripleta de $M'$? (Un tal $M'$ se llama \emph{emparejamiento perfecto}.)
\end{tcolorbox}

\begin{intuicion}
\DM\ es la versión ``en tres dimensiones'' del emparejamiento de un grafo
bipartito. Imaginamos $q$ trabajadores ($W$), $q$ tareas ($X$) y $q$ franjas
horarias ($Y$); cada tripleta de $M$ es una combinación viable
(trabajador, tarea, hora). Preguntamos si se pueden cubrir todos los
trabajadores, todas las tareas y todas las horas a la vez, sin repetir ninguno.
\end{intuicion}

\DM\ es uno de los problemas NP-completos originales de Karp (1972), y su
NP-completitud se demuestra por reducción desde 3-SAT. En el temario de la
asignatura aparece bajo el nombre \textbf{ACTRI} (acoplamiento de tripletas),
con la misma definición.

\subsection{Problema de la 4-Partición (\FPART)}

\begin{tcolorbox}[colback=orange!5, colframe=orange!40, title={\FPART\ (4-Partition)}]
\textbf{Entrada:} Un conjunto $A$ de $4m$ elementos, una cota $B \in \mathbb{Z}^+$,
y un tamaño $s(a) \in \mathbb{Z}^+$ por elemento, con
$\sum_{a \in A} s(a) = m B$ y $B/5 < s(a) < B/3$ para todo $a$.\\
\textbf{Pregunta:} ¿Puede particionarse $A$ en $m$ conjuntos disjuntos, cada uno
con suma $B$?
\end{tcolorbox}

La condición $B/5 < s(a) < B/3$ garantiza, por el mismo argumento de antes, que
cada conjunto de la partición tiene exactamente 4 elementos.

% ============================================================
\section{Demostración: \TPART\ es NP-completo}
% ============================================================

La demostración se organiza en tres partes:
\begin{enumerate}
  \item $\TPART \in \mathbf{NP}$.
  \item $\DM \leq_p \FPART$ (Teorema~\ref{thm:dm-to-4part}, basado en el
        Teorema 4.3 de Garey--Johnson).
  \item $\FPART \leq_p \TPART$ (Teorema~\ref{thm:4part-to-3part}, basado en el
        Teorema 4.4 de Garey--Johnson).
\end{enumerate}

\subsection{Parte I: \TPART\ $\in$ NP}

\begin{proposition}
\TPART\ pertenece a la clase NP.
\end{proposition}

\begin{proof}
Dado un certificado $\mathcal{C} = (S_1, S_2, \ldots, S_m)$
(una partición propuesta), el verificador comprueba en tiempo polinómico:
\begin{enumerate}
  \item Que los $S_i$ son disjuntos y cubren $A$: $O(n)$.
  \item Que $\sum_{a \in S_i} s(a) = B$ para todo $i$: $O(n)$.
\end{enumerate}
(El que cada $|S_i| = 3$ se deduce automáticamente de las cotas, pero también
puede comprobarse en $O(n)$.) Todo es verificable en tiempo $O(n)$ con
$n = 3m$. Por tanto, $\TPART \in \mathbf{NP}$.
\end{proof}

\subsection{Parte II: \DM\ $\leq_p$ \FPART\ (Teorema 4.3 de Garey--Johnson)}

\begin{theorem}[\DM\ se reduce a \FPART]
\label{thm:dm-to-4part}
\FPART\ es NP-completo en sentido fuerte; en particular, $\DM \leq_p \FPART$.
\end{theorem}

\begin{proof}
Sea $W = \{w_1,\ldots,w_q\}$, $X = \{x_1,\ldots,x_q\}$, $Y = \{y_1,\ldots,y_q\}$
y $M \subseteq W \times X \times Y$ una instancia arbitraria de \DM. Podemos
suponer sin pérdida de generalidad que $|M| \geq q$. Construiremos una instancia
de \FPART\ con $|A| = 4|M|$ elementos: cuatro por cada tripleta de $M$.

\medskip
\noindent\textbf{Elementos ``reales'' y ``ficticios''.}
Para cada $z \in W \cup X \cup Y$, sea $N(z)$ el número de tripletas de $M$ en las
que aparece $z$. Creamos $N(z)$ copias del elemento asociado a $z$: la copia
$z[1]$ se considera la \emph{real} (\textit{actual}), y las copias
$z[2], \ldots, z[N(z)]$ se consideran \emph{ficticias} (\textit{dummy}).

\begin{intuicion}
La idea de tener una copia ``real'' y varias ``ficticias'' por elemento resuelve
un problema de conteo. En \DM\ cada elemento de $W \cup X \cup Y$ debe usarse
\emph{exactamente una vez}; pero un mismo elemento puede aparecer en muchas
tripletas candidatas. Las copias ficticias absorben las apariciones que
\emph{no} se eligen en el emparejamiento, de modo que la copia real queda libre
para la tripleta ganadora. Los tamaños se diseñan para que ``real'' y ``ficticio''
encajen de formas distintas pero compatibles.
\end{intuicion}

\medskip
\noindent\textbf{Tamaños.}
Sea $r = 32q$. Los tamaños dependen de a qué conjunto ($W$, $X$ o $Y$) pertenece
$z$ y del índice del elemento dentro de su conjunto:
\[
\begin{aligned}
  s(w_i[1]) &= 10r^4 + ir + 1,      & 1 &\leq i \leq q, \\
  s(w_i[l]) &= 11r^4 + ir + 1,      & 1 &\leq i \leq q,\; 2 \leq l \leq N(w_i), \\[4pt]
  s(x_j[1]) &= 10r^4 + jr^2 + 2,    & 1 &\leq j \leq q, \\
  s(x_j[l]) &= 11r^4 + jr^2 + 2,    & 1 &\leq j \leq q,\; 2 \leq l \leq N(x_j), \\[4pt]
  s(y_k[1]) &= 10r^4 + kr^3 + 4,    & 1 &\leq k \leq q, \\
  s(y_k[l]) &= 8r^4 + kr^3 + 4,     & 1 &\leq k \leq q,\; 2 \leq l \leq N(y_k).
\end{aligned}
\]
A cada tripleta $m_l = (w_i, x_j, y_k) \in M$ le asociamos un elemento $u_l$ de
tamaño
\[
  s(u_l) = 10r^4 - kr^3 - jr^2 - ir + 8.
\]

\medskip
\noindent\textbf{Cota objetivo.}
Si sumamos $s(u_l)$ con los tamaños de los tres elementos correspondientes a
$w_i$, $x_j$, $y_k$, el total es $40r^4 + 15$ \emph{tanto si las tres copias son
reales como si las tres son ficticias} (la elección de los coeficientes lo
garantiza). Tomamos esa cantidad como cota\footnote{El libro de Garey y Johnson
imprime aquí $B = 40r^4 + 15\,|A|^3 + 15$; el término $15\,|A|^3$ es una errata:
cada 4-conjunto suma exactamente $40r^4 + 15$, que es el valor que debe tomar la
cota para que se cumpla $\sum_{a\in A} s(a) = |M|\,B_4$ (lo confirma el análisis
módulo $r$ de la propia demostración, que exige $B_4 \equiv 15 \pmod r$).}:
\[
  B_4 = 40r^4 + 15.
\]
Como $|A| = 4|M|$ y los elementos se reparten en grupos de cuatro, el número de
grupos objetivo es $m_4 = |A|/4 = |M|$.

\medskip
\noindent\textbf{Acotación.}
Se comprueba que cada tamaño está estrictamente entre $B_4/5$ y $B_4/3$, y que
$\sum s = m_4 \cdot B_4 = |M|\,(40r^4 + 15)$. Además, todos los tamaños están
acotados por
$12r^4 \leq 12 \cdot 8^4 \cdot |A|^4 < 2^{16} |A|^4$, es decir, \emph{polinómicamente}
en el número de elementos. Esta cota polinómica es lo que más adelante permitirá
invocar el Lema~\ref{lem:pseudo}.

\medskip
\noindent\textbf{Equivalencia de soluciones.}
\begin{itemize}
  \item ($\Rightarrow$) Si $M' \subseteq M$ es un emparejamiento perfecto,
        construimos una 4-partición así: por cada tripleta
        $m_l = (w_i,x_j,y_k) \in M'$ formamos el grupo
        $\{u_l, w_i[1], x_j[1], y_k[1]\}$ (las copias \emph{reales}); por cada
        tripleta $m_l \notin M'$ formamos $u_l$ junto con copias \emph{ficticias}.
        Hay suficientes copias ficticias para que esto siempre sea posible, y por
        la observación anterior cada grupo suma exactamente $B_4$.
  \item ($\Leftarrow$) Recíprocamente, sea una 4-partición. Analizando los tamaños
        módulo $r, r^2, r^3, r^4, r^5$ se demuestra que cada 4-conjunto debe
        contener un elemento de $W$, uno de $X$, uno de $Y$ y un $u_l$, y que los
        tres elementos de tipo $W,X,Y$ son o bien los tres reales o bien los tres
        ficticios. La colección de $3q$ elementos reales solo puede caber en $q$ de
        los 4-conjuntos; las tripletas de $M$ asociadas a esos $q$ conjuntos forman
        el emparejamiento perfecto buscado.
\end{itemize}

\begin{intuicion}
El truco de las potencias de $r$ (con $r = 32q$ ``suficientemente grande'') es una
forma de \emph{empaquetar} información en un solo número sin que las partes
interfieran. Cada potencia $r, r^2, r^3$ reserva una ``zona'' para codificar el
índice $i$, $j$, $k$ respectivamente; como nunca sumamos demasiados números, no
hay ``acarreo'' de una zona a otra. Es la misma idea que usa el temario de clase
con bloques de bits, solo que aquí las zonas se separan con potencias de $r$ en
lugar de potencias de $2$.
\end{intuicion}

La transformación es polinómica en $|M|$ y $q$, y genera números acotados
polinómicamente, luego es una transformación pseudopolinómica. Como \DM\ es
NP-completo (de hecho fuertemente, por ser sus números irrelevantes), el
Lema~\ref{lem:pseudo} da que \FPART\ es fuertemente NP-completo.
\end{proof}

\subsection{Parte III: \FPART\ $\leq_p$ \TPART\ (Teorema 4.4 de Garey--Johnson)}

\begin{theorem}[\FPART\ se reduce a \TPART]
\label{thm:4part-to-3part}
\TPART\ es NP-completo en sentido fuerte; en particular, $\FPART \leq_p \TPART$.
\end{theorem}

\begin{proof}
Para que la transformación sea \emph{pseudopolinómica} no reducimos el \FPART\
general, sino el subproblema en el que $\max\{s(a): a \in A\} \leq 2^{16} |A|^4$
(que es exactamente el que produce el Teorema~\ref{thm:dm-to-4part}, y que sigue
siendo fuertemente NP-completo). Sea entonces
$A = \{a_1, \ldots, a_{4n}\}$, cota $B$, tamaños $s(a)$ con
$B/5 < s(a) < B/3$ y $s(a) \leq 2^{16} |A|^4$.

\medskip
\noindent\textbf{Construcción.}
La instancia de \TPART\ tendrá $24n^2 - 3n$ elementos, de tres tipos:

\begin{description}
  \item[Regulares.] Por cada elemento $a_i \in A$, un elemento ``regular'' $w_i$ con
  \[
    s'(w_i) = 4\,(5B + s(a_i)) + 1.
  \]
  \item[De emparejamiento.] Por cada \emph{par} $a_i, a_j \in A$ (con $i<j$), dos
  elementos ``de emparejamiento'' $u[i,j]$ y $\bar u[i,j]$ con
  \[
    s'(u[i,j]) = 4\,(6B - s(a_i) - s(a_j)) + 2,
    \qquad
    s'(\bar u[i,j]) = 4\,(5B + s(a_i) + s(a_j)) + 2.
  \]
  \item[De relleno.] Para $1 \leq k \leq 8n^2 - 3n$, un elemento ``de relleno''
  $u_k^*$ con
  \[
    s'(u_k^*) = 20B.
  \]
\end{description}
La cota de la instancia de \TPART\ es
\[
  B' = 64B + 4.
\]

\begin{intuicion}
La clave de toda esta construcción es el \textbf{resto módulo 4} de cada tamaño:
\begin{center}
\begin{tabular}{lll}
\toprule
Tipo & Forma del tamaño & $s' \bmod 4$ \\
\midrule
Regular        & $4(\cdots)+1$ & $1$ \\
Emparejamiento & $4(\cdots)+2$ & $2$ \\
Relleno        & $20B = 4\cdot 5B$ & $0$ \\
\bottomrule
\end{tabular}
\end{center}
Como la cota es $B' = 64B+4 \equiv 0 \pmod 4$, cualquier trío que sume $B'$ debe
tener restos que sumen $0 \pmod 4$. Con restos disponibles $\{1, 2, 0\}$, las
combinaciones de tres que dan $0 \pmod 4$ son:
\[
  0+0+0 \equiv 0, \qquad 1+1+2 = 4 \equiv 0, \qquad 2+2+0 = 4 \equiv 0.
\]
La primera, $0+0+0$ (tres rellenos), queda descartada por el \emph{valor} de la
suma: tres rellenos suman $3 \cdot 20B = 60B \neq 64B + 4 = B'$. Es decir, el
criterio módulo $4$ es necesario pero no suficiente, y hay que combinarlo con el
tamaño objetivo. Quedan así solo dos tipos de trío válido: o bien
\{regular, regular, emparejamiento\}, o bien
\{emparejamiento, emparejamiento, relleno\}. Esta rigidez es lo que hace que la
3-partición ``imite'' a la 4-partición original.
\end{intuicion}

\medskip
\noindent\textbf{Equivalencia de soluciones.}
\begin{itemize}
  \item ($\Rightarrow$) Si $\{a_i, a_j, a_k, a_l\}$ es un 4-conjunto de la
        4-partición (con suma $B$), lo dividimos en dos pares, digamos
        $\{a_i, a_j\}$ y $\{a_k, a_l\}$, y formamos los dos tríos
        \[
          \{w_i, w_j, u[i,j]\}
          \quad\text{y}\quad
          \{w_k, w_l, \bar u[i,j]\}.
        \]
        Cada uno suma exactamente $B'$: por ejemplo el primero suma
        $4(5B+s(a_i)) + 1 + 4(5B+s(a_j)) + 1 + 4(6B - s(a_i) - s(a_j)) + 2
        = 4(16B) + 4 = 64B + 4 = B'$.
        Repitiendo para los $n$ 4-conjuntos obtenemos $2n$ tríos que usan todos los
        regulares y $n$ pares emparejados, dejando $8n^2 - 3n$ pares de
        emparejamiento y $8n^2 - 3n$ elementos de relleno. Como la suma de un par emparejado
        $u[i,j] + \bar u[i,j] = 44B + 4 = B' - 20B$, cada par sobrante se agrupa con
        un relleno para completar la 3-partición.
  \item ($\Leftarrow$) Si tenemos una 3-partición, el análisis módulo 4 (la caja de
        intuición) dice que cada trío es de uno de los dos tipos permitidos.
        Mediante un argumento de intercambio (los elementos de emparejamiento de un
        mismo par tienen tamaños que ``encajan'' entre sí) se reorganizan los tríos
        hasta que cada relleno aparece junto a un par emparejado completo
        $u[i,j], \bar u[i,j]$. Eso obliga a que los tríos de tipo
        \{regular, regular, emparejamiento\} agrupen los regulares en $n$ pares de
        tríos cuyos cuatro regulares suman $84B + 4$, lo que equivale a que los
        cuatro elementos originales de $A$ sumen $B$. Esos $n$ grupos de cuatro
        forman la 4-partición.
\end{itemize}

\medskip
\noindent\textbf{Acotación y carácter pseudopolinómico.}
Cada tamaño está estrictamente entre $B'/4 = 16B+1$ y $B'/2 = 32B+2$, y
$\sum s' = (8n^2 - n) B'$. Como en $A$ los tamaños estaban acotados por
$2^{16}|A|^4$, los tamaños de la instancia de \TPART\ también quedan acotados por
un polinomio en $|A|$. Por tanto la transformación es pseudopolinómica y, por el
Lema~\ref{lem:pseudo}, \TPART\ es fuertemente NP-completo.
\end{proof}

\begin{remark}[Por qué la restricción del subproblema es necesaria]
Garey y Johnson advierten un punto sutil: si se intentara aplicar esta misma
construcción al \FPART\ \emph{general} (sin acotar $\max\{s(a)\}$), la
transformación seguiría siendo correcta como reducción polinómica ordinaria, pero
\textbf{no} bastaría para la NP-completitud \emph{fuerte}, porque los números
podrían crecer demasiado. Restringir al subproblema con $\max\{s(a)\}$ acotado
polinómicamente —que es justo lo que entrega el paso anterior— es lo que mantiene
la dureza fuerte. Este detalle es la razón técnica de fondo por la que la cadena se
diseña con cuidado en cada eslabón.
\end{remark}

\subsection{Resultado Principal}

\begin{theorem}[NP-completitud de \TPART]
El problema de la 3-Partición es NP-completo (de hecho, fuertemente NP-completo).
\end{theorem}

\begin{proof}
Se sigue de lo anterior:
\begin{enumerate}
  \item $\TPART \in \mathbf{NP}$ (verificación en tiempo lineal).
  \item \DM\ es fuertemente NP-completo (Karp, 1972; sus números son irrelevantes).
  \item $\DM \leq_p \FPART$ por transformación pseudopolinómica
        (Teorema~\ref{thm:dm-to-4part}).
  \item $\FPART \leq_p \TPART$ por transformación pseudopolinómica
        (Teorema~\ref{thm:4part-to-3part}).
\end{enumerate}
Por transitividad, \TPART\ es NP-difícil; combinado con $\TPART \in \mathbf{NP}$,
es NP-completo. Como ambas reducciones son pseudopolinómicas y mantienen los
números acotados polinómicamente, el Lema~\ref{lem:pseudo} da además que es
fuertemente NP-completo.
\end{proof}

% ============================================================
\section{¿Por qué dos reducciones y no una?}
\label{sec:dos-reducciones}
% ============================================================

Una duda razonable al ver la cadena $\DM \leq_p \FPART \leq_p \TPART$ es por qué
se encadenan \emph{dos} reducciones en lugar de hacer una sola.

Para probar que \TPART\ es NP-difícil bastaría, en principio, \emph{una} reducción
desde \emph{cualquier} problema NP-completo. Encadenar dos pasos \textbf{no es una
obligación lógica}, sino una conveniencia técnica: construir directamente
$\DM \leq_p \TPART$ sería muy difícil de escribir y de verificar, mientras que el
escalón intermedio \FPART\ descompone esa dificultad en dos saltos manejables.
Como $\leq_p$ es transitiva, las dos reducciones encadenadas equivalen a una sola
$\DM \leq_p \TPART$.

\begin{intuicion}
Por ende, el hecho de hacerlo así es por comodidad, no necesidad: la demostración podría hacerse
con una única reducción, pero saldría intratable de escribir y de comprobar. El
bloque intermedio \FPART\ existe precisamente para que cada salto sea manejable.
\end{intuicion}

% ============================================================
\section{Comparación: la cadena del libro frente a la del temario}
\label{sec:comparacion}
% ============================================================

El temario de la asignatura (Tema 4, Serafín Moral) demuestra un resultado
\emph{emparentado pero distinto} al de este trabajo. Conviene tenerlo claro porque
los dos usan tripletas y se parecen superficialmente.

\begin{center}
\begin{tabular}{p{4.3cm} p{5.0cm} p{4.6cm}}
\toprule
 & \textbf{Temario (Serafín Moral)} & \textbf{Este trabajo (Garey--Johnson)} \\
\midrule
Problema final & \PART\ clásica (en \textbf{2} grupos) & \TPART\ (en grupos de \textbf{3}) \\
Cadena & 3-SAT $\to$ ACTRI $\to$ SUMA $\to$ \PART & $\DM \to \FPART \to \TPART$ \\
Tipo de dureza & NP-completo (\textbf{débil}) & NP-completo (\textbf{fuerte}) \\
¿Admite alg.\ pseudopolinómico? & Sí & No (salvo $\mathbf{P}=\mathbf{NP}$) \\
Dificultad técnica & Media & Alta \\
\bottomrule
\end{tabular}
\end{center}

\begin{remark}
Los dos enfoques son \textbf{correctos}; simplemente prueban teoremas diferentes.
El del temario llega a la Partición en \emph{dos} grupos, que es NP-completo
\emph{débil} (de hecho tiene algoritmo pseudopolinómico por programación dinámica).
El de este trabajo llega a la 3-Partición en grupos de \emph{tres}, que es
fuertemente NP-completo: un resultado más fino y técnicamente más exigente. Por eso
el temario puede permitirse una codificación binaria por zonas (con potencias de
$2$) mientras que Garey--Johnson debe cuidar que todos los números queden acotados
polinómicamente.
\end{remark}

\paragraph{El parentesco.} Pese a las diferencias, las dos cadenas comparten puntos:
ambas parten de un problema lógico (3-SAT), pasan por un problema de tripletas
(ACTRI~$=$~\DM) y terminan en un problema de partición de números. La técnica común
es la \emph{codificación por zonas}: empaquetar varios índices en un solo número de
forma que no se ``mezclen'' al sumar (con potencias de $2$ en el temario, con
potencias de $r=32q$ en Garey--Johnson).

% ============================================================
\section{Relación con las Técnicas de Reducción del Temario}
% ============================================================

La demostración emplea, en distinto grado, las técnicas de reducción del Tema 4:

\begin{description}
  \item[Reemplazamiento Local con Refuerzo.]
        En $\DM \leq_p \FPART$ cada tripleta de \DM\ se transforma en un elemento
        $u_l$, y cada elemento de $W\cup X\cup Y$ en copias real/ficticias. Son los
        ``elementos adicionales'' (las copias ficticias y los $u_l$) los que
        \emph{refuerzan} la construcción para forzar la equivalencia, exactamente la
        idea de esta técnica.

  \item[Reemplazamiento Local con Refuerzo (de nuevo).]
        En $\FPART \leq_p \TPART$ los elementos de relleno $u_k^*$ y los de
        emparejamiento $u[i,j], \bar u[i,j]$ son refuerzos añadidos para que la
        3-partición solo pueda imitar a la 4-partición original.

  \item[Diseño de Componentes.]
        El uso de tipos de elemento con distinto resto módulo 4 (regular,
        emparejamiento, relleno), que se ``conectan'' obligatoriamente en tríos de
        dos formas posibles, es un ejemplo de diseño de componentes: distintas
        piezas con un comportamiento prefijado que encajan de manera controlada.
\end{description}

% ============================================================
\section{Ejemplo Numérico (cadena del temario: ACTRI $\to$ SUMA)}
% ============================================================

\begin{remark}
El siguiente ejemplo ilustra la \emph{codificación por zonas} sobre la cadena del
temario ($\text{ACTRI} \to \text{SUMA}$), no sobre la reducción
$\DM \to \FPART$ de Garey--Johnson (cuyos tamaños, con potencias de $r=32q$, son
demasiado grandes para tabularse a mano). Se incluye porque muestra de forma
transparente la \emph{misma idea} de empaquetar índices en zonas que está detrás de
ambas cadenas.
\end{remark}

\begin{example}[Codificación por zonas]
Sea la instancia de ACTRI ($=\DM$) con
\[
  W = \{w_1, w_2, w_3\},\quad X = \{x_1, x_2, x_3\},\quad Y = \{y_1, y_2, y_3\},
\]
\[
  M = \{(w_2,x_1,y_3),\, (w_2,x_2,y_2),\, (w_3,x_2,y_1),\,
        (w_3,x_3,y_2),\, (w_1,x_3,y_2)\}.
\]
Hay $q=3$, $k = |M| = 5$, y se toma $p = \lfloor \log_2 k \rfloor + 1 = 3$ bits por
zona. Cada tripleta se codifica como un número con nueve zonas de 3 bits (una por
cada elemento de $W \cup X \cup Y$), poniendo un $1$ en las zonas de sus tres
elementos:

\smallskip
\begin{center}
\begin{tabular}{cccccccccc}
\toprule
Tripleta & $w_1$ & $w_2$ & $w_3$ & $x_1$ & $x_2$ & $x_3$
         & $y_1$ & $y_2$ & $y_3$ \\
\midrule
$(w_2,x_1,y_3)$ & 000 & 001 & 000 & 001 & 000 & 000 & 000 & 000 & 001 \\
$(w_2,x_2,y_2)$ & 000 & 001 & 000 & 000 & 001 & 000 & 000 & 001 & 000 \\
$(w_3,x_2,y_1)$ & 000 & 000 & 001 & 000 & 001 & 000 & 001 & 000 & 000 \\
$(w_3,x_3,y_2)$ & 000 & 000 & 001 & 000 & 000 & 001 & 000 & 001 & 000 \\
$(w_1,x_3,y_2)$ & 001 & 000 & 000 & 000 & 000 & 001 & 000 & 001 & 000 \\
\midrule
$B$             & 001 & 001 & 001 & 001 & 001 & 001 & 001 & 001 & 001 \\
\bottomrule
\end{tabular}
\end{center}

\smallskip\noindent
La solución de ACTRI es $M' = \{(w_2,x_1,y_3),\,(w_3,x_2,y_1),\,(w_1,x_3,y_2)\}$
(tripletas 1, 3 y 5). Al sumar sus tres números se obtiene exactamente $B$ (un $1$
en cada zona, sin acarreos, porque $p=3$ bits bastan para que sumar a lo sumo $k=5$
unos nunca desborde una zona). Esto corresponde, en la cadena del temario, a una
solución del problema SUMA equivalente.
\end{example}

% ============================================================
\section{Aplicaciones y Relevancia}
% ============================================================

La NP-completitud fuerte de \TPART\ se usa para demostrar la intratabilidad de
muchos problemas, porque permite reducir desde un problema cuya dureza no depende de
que los números sean grandes:

\begin{itemize}
  \item \textbf{Planificación multiprocesador y secuenciación:} numerosos problemas
        de \emph{scheduling} en máquinas idénticas son NP-completos en sentido fuerte
        gracias a reducciones desde \TPART. Este fue el contexto original de
        Garey y Johnson.

  \item \textbf{Empaquetamiento (bin packing, empaquetado de rectángulos):} la
        NP-dureza de varias variantes se obtiene por reducción desde \TPART.

  \item \textbf{Bandwidth de grafos:} en el propio apéndice de Garey--Johnson, los
        problemas \textsc{Bandwidth} y \textsc{Directed Bandwidth} [GT40, GT41] se
        marcan como NP-completos por transformación desde 3-Partition.

  \item \textbf{Juegos y rompecabezas:} trabajos posteriores
        (p.\,ej. Demaine y coautores sobre Tetris) usan reducciones desde \TPART\ o
        variantes para probar NP-dureza.
\end{itemize}

% ============================================================
\section{Conclusión}
% ============================================================

Se ha demostrado que el \textbf{Problema de la 3-Partición} es
\textbf{fuertemente NP-completo}, siguiendo fielmente la construcción de
Garey y Johnson:
\[
  \underbrace{\DM}_{\text{NP-completo}} \;\leq_p\;
  \underbrace{\FPART}_{\text{fuert.\ NP-completo}} \;\leq_p\;
  \underbrace{\TPART}_{\text{fuert.\ NP-completo}}.
\]

Los puntos clave son:
\begin{enumerate}
  \item La verificación de una solución es polinómica ($\TPART \in \mathbf{NP}$).
  \item En $\DM \leq_p \FPART$, las copias real/ficticias y los elementos $u_l$,
        junto con la codificación por potencias de $r = 32q$, hacen corresponder
        emparejamientos con 4-particiones manteniendo los números acotados
        polinómicamente.
  \item En $\FPART \leq_p \TPART$, los tres tipos de elemento (regular,
        emparejamiento, relleno) y el análisis módulo 4 fuerzan que toda 3-partición
        imite a la 4-partición original.
  \item Ambas reducciones son \emph{pseudopolinómicas}, por lo que el Lema 4.1
        propaga la dureza fuerte hasta \TPART.
\end{enumerate}

La 3-Partición es un ejemplo paradigmático: su enunciado es simple, pero su
intratabilidad es ``robusta'' (no depende de números grandes), y por eso es una de
las herramientas más empleadas para demostrar NP-dureza fuerte de otros problemas.

% ============================================================
\begin{thebibliography}{99}
\small

\bibitem{garey1979}
M.\,R. Garey y D.\,S. Johnson.
\textit{Computers and Intractability: A Guide to the Theory of NP-Completeness}.
W.\,H. Freeman, San Francisco, 1979. Teoremas 4.3 y 4.4 (pp.~96--101);
problema [SP15].

\bibitem{karp1972}
R.\,M. Karp.
``Reducibility among combinatorial problems.''
En R.\,E. Miller y J.\,W. Thatcher (eds.),
\textit{Complexity of Computer Computations}, Plenum Press, 1972, pp.~85--103.

\bibitem{moral}
S. Moral.
\textit{Tema 4: NP-completitud.} Apuntes de Complejidad Computacional,
Universidad de Granada, 2024.

\end{thebibliography}

\end{document}
